\begin{textsample}{NEG Dim 4 – 2014\_c5 – Score 1.00 – 2014\_c5/enade\_2014\_art\_vis\_fab093e004.txt}  \label{ex:f4_neg_204}
CAPÍTULO 5 DISTRIBUIÇÃO DOS CONCEITOS 5.1 PANORAMA NACIONAL DA DISTRIBUIÇÃO DOS CONCEITOS O Conceito Enade é um indicador de qualidade que avalia o desempenho dos estudantes a partir dos resultados do ENADE.
Os valores possíveis vão de 1 ( pior situação ) a 5 ( melhor situação ).
As Unidades de Observação com apenas um ou sem nenhum concluinte participante não obtêm o Conceito Enade, ficando Sem Conceito ( SC ).
A Tabela 5.1 apresenta a quantidade e distribuição de cursos de Artes Visuais ( Licenciatura ) participantes do ENADE/2014, por faixa de conceito e Grande Região.
Mais uma vez, ressaltando que a diferença entre os cursos tabulados neste capítulo e no capítulo 2 são os cursos sem conceito, em princípio, aqueles sem alunos concluintes que participassem da prova ou que tivessem tão somente um aluno concluinte²².
Observando -se os dados da Tabela 5.1, nota -se que, dos 108 cursos participantes, 47 ( 43,5 \% ) classificaram -se com conceito 3, o valor modal.
Este foi também o conceito modal em mais quatro regiões : Norte ( 62,5 \% ), Nordeste ( 36,4 \% ), Sudeste ( 53,8 \% ) e Sul ( 33,3 \% ).
Na região Centro-Oeste, a moda foi o conceito 4 ( 66,7 \% ).
O conceito 4 foi o segundo mais frequente em nível nacional ( 21,3 \%, correspondendo a 23 cursos ) e o conceito 2, o terceiro ( 17,6 \%, correspondendo a 19 cursos ).
Houve, ainda, nove cursos ( 8,3 \% ) que receberam conceito 5 e quatro cursos ( 3,7 \% ) que receberam conceito 1.
Dos 108 cursos de Artes Visuais ( Licenciatura ), apenas seis cursos ( 5,6 \% ) ficaram sem conceito ( SC ). ²² Estes cursos com somente um concluinte não foram considerados no capítulo 2.
A região Norte participou com oito cursos ou 7,4 \% do total nacional.
O conceito 3 foi atribuído a cinco cursos dos oito participantes, o que equivale a 62,5 \% do total regional, além de ser o conceito modal na região.
Nenhum curso ficou sem conceito ( SC ) ou com conceitos 1 ou 5.
Dois cursos foram avaliados com conceito 2 e um curso, com conceito 4, correspondendo, respectivamente, a 25,0 \% e a 12,5 \% do total da região.
A região Nordeste participou com 22 cursos ou 20,4 \% do total nacional.
Nesta região, oito cursos, 36,4 \% em termos regionais, obtiveram conceito 3, o conceito modal para a região.
Os conceitos 2 e 4 foram atribuídos a, respectivamente, quatro e a cinco cursos ( 18,2 \% e 22,7 \% ).
Dois cursos ficaram sem conceito ( SC ), um curso ( 4,5 \% ) recebeu o conceito 1 e os dois cursos restantes ( 9,1 \% ) receberam o conceito 5, nessa região.
Dos 39 cursos participantes da região Sudeste ( 36,1 \% do total nacional ), o conceito 3 também foi o conceito modal, tendo sido atribuído a 21 cursos ( 53,8 \% ) da região.
A nenhum curso da região foi atribuído conceito 1.
Seis cursos ( 15,4 \% ) receberam conceito 2, quatro cursos ( 10,3 \% ) obtiveram conceito 4 e a sete cursos ( 17,9 \% ) foi atribuído conceito 5.
Um curso ( 2,6 \% ) ficou sem conceito ( SC ).
Dos 33 cursos da região Sul, três ficaram sem conceito ( SC ) e os demais ficaram distribuídos em quatro das cinco faixas de conceitos.
Há predominância do conceito 3 com 33,3 \% ( conceito modal ), ou seja, 11 dos 33 cursos participantes na região Sul.
O conceito 2 foi atribuído a sete cursos ( 21,2 \% ) e o conceito 4, a nove cursos ( 27,3 \% ).
A nenhum curso da região foi atribuído conceito 5.
Já dos seis cursos participantes na região Centro-Oeste, dois ( 33,3 \% ) receberam conceito 3 e quatro ( 66,7 \% ), o conceito modal, 4.
Nesta região, nenhum curso ficou sem conceito ( SC ) ou com conceito 1, 2 ou 5.
O Gráfico 5.1 apresenta a distribuição do Conceito Enade segundo Grande Região : poligonais mais à esquerda representam Grandes Regiões com pior distribuição, e poligonais mais à direita, Grandes Regiões com melhores conceitos.
A região Centro-Oeste ( linha preta ) apresenta os cursos com conceitos mais elevados, comparativamente aos cursos oferecidos nas demais regiões – regiões Norte ( linha azul ), Nordeste ( linha verde ), Sudeste ( linha vermelha ) e Sul ( linha roxa ) – que apresentam as distribuições com valores relativamente mais baixos. 5.2 CONCEITOS POR CATEGORIA ADMINISTRATIVA E POR GRANDE REGIÃO A Tabela 5.2 apresenta a distribuição dos cursos participantes do ENADE/2014 de Artes Visuais ( Licenciatura ), por Categoria Administrativa, de acordo com os conceitos por eles alcançados, segundo as Grandes Regiões brasileiras.
Dos 108 cursos participantes, 56 ( 51,9 \% ) eram ministrados em Instituições Públicas e 52 ( 48,1 \% ), em Instituições Privadas.
De acordo com as informações da Tabela 5.2, em termos nacionais, entre Instituições Públicas e Privadas : das nove IES cujos cursos receberam conceito 5, cinco eram Públicas e as outras quatro, Privadas.
Dos 56 cursos participantes de IES Públicas, o conceito 3 foi o valor modal, atribuído a 21 cursos ( 37,5 \% da categoria ).
Entre os demais cursos de Instituições Públicas participantes, um único curso ( 1,8 \% ) foi avaliado com conceito 1, oito cursos ( 14,3 \% ) obtiveram conceito 2, 16 cursos ( 28,6 \% ) receberam conceito 4 e, como já comentado, cinco cursos ( 8,9 \% ) foram \textbf{alocados} ao conceito 5.
Nesta categoria, cinco cursos ficaram sem conceito ( SC ).
Na rede privada, o conceito modal também foi 3, com 26 cursos dos 52 da categoria.
Entre os demais cursos de IES Privadas participantes, três receberam conceito 1, e onze, conceito 2.
O conceito 4 foi atribuído a sete cursos e o conceito 5, a quatro cursos.
Nesta Categoria Administrativa, apenas um curso ficou sem conceito ( SC ).
Na análise por região, observa -se que, na região Norte, as Instituições Privadas participaram com dois cursos ( 25,0 \% do total regional ), ambos \textbf{alocados} ao conceito 3.
Nenhum curso de Artes Visuais ( Licenciatura ) de IES Privada situada nesta região ficou sem conceito ( SC ).
As Instituições Públicas participaram com seis cursos na região Norte ( 60,0 \% do total regional ), dos quais três obtiveram conceitos 4, o valor modal, seguido por dois cursos ( 33,3 \% ) que receberam conceito 2 e um único curso que obteve conceito 4 ( 16,7 \% ).
Nenhum curso ficou sem conceito ( SC ) ou recebeu conceito 1 nesta combinação de categoria e região.
Na região Nordeste, a rede privada concentrou dois dos 22 cursos participantes, o equivalente a 9,1 \% do total da região.
Um desses cursos ficou \textbf{alocado} ao conceito 1, e o outro, ao conceito 5.
As Instituições Públicas dessa região participaram com 20 cursos ( 90,9 \% ), a maior das proporções desta categoria em uma região, dos quais oito ( 40,0 \% ) obtiveram conceito 3, o valor modal.
Os demais cursos receberam os conceitos : 2 ( quatro cursos ), 4 ( cinco cursos ) e 5 ( um curso ).
Nesta combinação de categoria e região, dois cursos ficaram sem conceito ( SC ), e nenhum curso recebeu conceito 1.
Na região Sudeste, a proporção de cursos da rede privada, 71,8 \%, foi mais elevada do que nas demais regiões brasileiras, correspondendo a 28 dos 39 cursos participantes.
O conceito modal em IES Privadas na região Sudeste foi 3, atribuído a 15 cursos.
Nesta combinação de Categoria Administrativa e Grande Região, um curso não recebeu conceito ( SC ) e nenhum curso recebeu conceito 1.
Os demais cursos foram avaliados com os conceitos 2 ( seis cursos ), 4 ( três cursos ) e 5 ( três cursos ).
Dos 11 cursos oferecidos por Instituições Públicas na região Sudeste, seis receberam conceito 3 ( conceito modal ).
Aos demais cursos foram atribuídos os conceitos 4 ( um curso ) e 5 ( quatro cursos ).
Nenhum curso ficou sem conceito ( SC ) ou recebeu conceitos 1 ou 2.
As Instituições Privadas concentraram 17 dos 33 cursos participantes da região Sul, 51,5 \% do total regional.
Desses, sete ficaram com conceito 3, o conceito modal.
Nesta combinação de Categoria Administrativa e Grande Região, dois cursos ficaram com conceito 1, cinco cursos obtiveram conceito 2 e três cursos receberam conceito 4.
Nenhum curso ficou sem conceito ( SC ) ou recebeu conceito 5.
As Instituições Públicas na região Sul participaram com 16 cursos ( 48,5 \% do total da região ).
O conceito modal foi 4, atribuído a seis cursos.
Os demais cursos receberam os conceitos 1 ( um curso ), 2 ( dois cursos ) e 3 ( quatro cursos ).
Três cursos ficaram sem conceito ( SC ), e a nenhum curso de IES Publica situado nessa região foi atribuído conceito 5.
Na região Centro-Oeste, os cursos participantes distribuíram -se de modo equilibrado entre as duas Categorias Administrativas : três estavam vinculados a IES privadas e outros três, a IES Públicas.
Independentemente da categoria em que se enquadrasse, nenhum dos seis cursos participantes ficou sem conceito ( SC ) ou recebeu os conceitos 1, 2 ou 5.
Dos três cursos vinculados a Instituições Privadas, dois receberam conceito 3 e o outro, conceito 4.
Quanto aos três cursos oferecidos por Instituições Públicas, todos foram avaliados com conceito 4.
O Gráfico 5.2 apresenta a distribuição do Conceito Enade segundo a Categoria Administrativa da IES.
Os cursos em IES Públicas ( linha azul ) apresentam uma poligonal mais à direita do que os cursos em IES Privadas ( linha verde ) e, consequentemente, uma distribuição de Conceitos Enade com valores maiores. 5.3 CONCEITOS POR ORGANIZAÇÃO ACADÊMICA E POR GRANDE REGIÃO Na Tabela 5.3, encontra -se a distribuição dos conceitos atribuídos aos cursos participantes do ENADE/2014 na Área de Artes Visuais ( Licenciatura ), por Organização Acadêmica, segundo as Grandes Regiões brasileiras.
Dos 108 cursos de Artes Visuais ( Licenciatura ) participantes, 77 eram oferecidos em Universidades, 11 em Centros Universitários e 20 em Faculdades.
Esta distribuição corresponde a, respectivamente, 71,3 \%, 10,2 \% e 18,5 \% dos cursos.
De acordo com os dados apresentados, de todos os nove cursos avaliados com conceito 5, seis eram vinculados a Universidades.
Seis cursos oferecidos em Universidades ficaram sem conceito ( SC ).
Esse tipo de Organização Acadêmica teve o conceito 3 como modal, com 34 cursos ( 44,2 \% ).
Os demais cursos de Universidades avaliados receberam os conceitos 1 ( três cursos ), 2 ( dez cursos ), 4 ( 18 cursos ) e conceito 5 ( seis cursos, como já mencionado ).
Entre os cursos em Centros Universitários, o conceito modal foi 3, atribuído a cinco cursos ( 45,5 \% ).
Os outros cursos receberam os conceitos 2 ( três cursos ), 4 ( dois cursos ) e 5 ( um curso ).
Neste tipo de Organização Acadêmica, nenhum curso ficou sem conceito ( SC ) ou com conceito 1.
Nas Faculdades, dos 20 cursos participantes, oito ( 40,0 \% ) ficaram \textbf{alocados} no conceito 3, conceito modal.
Os demais cursos obtiveram conceitos 1 ( um curso ), 2 ( seis cursos ), 4 ( três cursos ) e 5 ( dois cursos ).
Nenhum curso ficou sem conceito ( SC ) neste tipo de Organização Acadêmica.
Considerando -se separadamente as regiões brasileiras, verifica -se que, na região Norte, as Universidades concentraram sete dos oito cursos participantes.
Quatro cursos ficaram \textbf{alocados} ao conceito 3, conceito modal, seguidos de dois cursos que obtiveram conceito 2, enquanto somente um curso recebeu conceito 4.
Nenhum curso ficou sem conceito ( SC ) ou recebeu os demais conceitos nesta categoria e região.
Na região Norte, nenhum Centro Universitário ofereceu curso na área de Artes Visuais ( Licenciatura ), e apenas uma Faculdade da região ofereceu o curso, ao qual foi atribuído conceito 3.
Na região Nordeste, as Universidades participaram com 21 dos 22 cursos na área de Artes Visuais ( Licenciatura ).
O conceito modal foi 3, atribuído a oito cursos.
Os demais cursos receberam os conceitos 1 ( um curso ), 2 ( quatro cursos ), 4 ( cinco cursos ) e 5 ( um curso ).
Dois cursos oferecidos por Universidades no Nordeste ficaram sem conceito ( SC ).
Nenhum Centro Universitário situado na região Nordeste ofereceu curso na área de Artes Visuais ( Licenciatura ), e apenas uma Faculdade da região ofereceu o curso, ao qual foi atribuído conceito 5.
Na região Sudeste, as Universidades concentraram 21 dos 39 cursos participantes.
Entre os cursos vinculados a Universidades, nessa região, o conceito modal foi 3, atribuído a 13 cursos.
Aos demais cursos, foram atribuídos os conceitos 2 ( um curso ), 4 ( um curso ) e 5 ( cinco cursos ).
Apenas um curso ficou sem conceito ( SC ), e a nenhum foi atribuído conceito 1.
Os Centros Universitários participaram com 8 cursos na região Sudeste, metade ( quatro ) dos quais obteve o conceito 3 ( conceito modal ).
Nessa região, nenhum curso oferecido por Centro Universitário ficou sem conceito ( SC ) ou recebeu conceito 1.
Os demais cursos se distribuíram nos conceitos 2 ( dois cursos ), 4 ( um curso ) e 5 ( um curso ).
As Faculdades foram representadas por dez cursos na região Sudeste dos quais quatro receberam o conceito modal, 3.
Os demais receberam os conceitos 2 ( três cursos ), 4 ( dois cursos ) e 5 ( um curso ).
Nenhum curso vinculado a esse tipo de organização acadêmica ficou sem conceito ( SC ) ou recebeu conceito 1 nessa região.
Dos 33 cursos da região Sul, 25 eram de Universidades, para os quais o conceito modal foi 3, atribuído a nove cursos, seguido de mais oito cursos que receberam o conceito 4.
Nesse tipo de organização, nenhum curso recebeu conceito 5.
Os demais cursos receberam os conceitos 1 ( dois cursos ) e 2 ( três cursos ).
Três cursos oferecidos em Universidades nessa região ficaram sem conceito ( SC ).
Os dois cursos participantes de Centros Universitários da região Sul distribuíram -se nos conceitos 2 e 3 ( um curso cada ).
Já os seis cursos vinculados a Faculdades na região Sul receberam os conceitos 1 ( um curso ), 2 ( três cursos, Na região Centro-Oeste, três dos seis cursos participantes eram oferecidos em Universidades, todos avaliados com conceito 4.
Nessa região, independentemente do tipo de organização acadêmica, nenhum curso ficou sem conceito ( SC ).
Apenas um Centro Universitário ofereceu, na região Centro-Oeste, o curso de Artes Visuais ( Licenciatura ), ao qual foi atribuído conceito 4.
As Faculdades da região, por sua vez, participaram com dois cursos, ambos avaliados com conceito 3.
O Gráfico 5.3 apresenta a distribuição do Conceito Enade segundo a Organização Acadêmica da IES.
Os cursos em Universidades ( linha azul ) apresentam uma poligonal levemente mais à direita do que as demais, denotando uma melhor distribuição de Conceitos Enade.
No outro extremo, temos os cursos oferecidos em Faculdades ( linha vermelha ) com a poligonal ligeiramente mais à esquerda e, consequentemente, uma distribuição de Conceitos Enade com menores valores.
IES de Centros Universitários ( linha verde ) apresentam uma situação intermediária.

% matched lemmas: alocar
\end{textsample}
